\documentclass[10pt]{beamer}

\usetheme{m}

\usepackage{booktabs}
\usepackage[scale=2]{ccicons}

%\usepackage{pgfplots}
%\usepgfplotslibrary{dateplot}

\usepackage{listings}


\title{}
\subtitle{Intel Threading Building Blocks}
\date{}
\author{Luis María Costero Valero\\Jesús Javier Domenech Arellano\\Hristo
  Ivanov Ivanov}
\institute{23 Noviembre 2015}

\titlegraphic{\includegraphics[scale=0.4]{bird-illustration-tbb.png}}
\begin{document}

\maketitle





%%======= INDICE ========================================
%\begin{frame}
%  \frametitle{Índice}
%  \setbeamertemplate{section in toc}[sections numbered]
%  \tableofcontents[hideallsubsections]
%\end{frame}

%\section{Introduction} %=== Esrto genera una página de sección.

\begin{frame}
  \frametitle{Intel Threading Building Blocks}
  \emph{ \
    Es una librería \texttt{C++} desarrollada por \emph{Intel}, cuyo objetivo
    es facilitar la creación de programas paralelos.
    La librería ofrece diferentes estructuras de datos y algoritmos que
    ofrecen al programador una abstracción de alto nivel sobre la máquina.
    \begin{itemize}
      \item[] Fabricante: Intel.
      \item[] Compatibilidad: FreeBSD, Linux, Solaris, OS X, Windows, Android.
      \item[] Versión \texttt{1.0} el 29 de Agosto de 2006.
      \item[] Versión actual: 4.4 / 25 de Agosto de, 2015.
      \item[] Licencia: dual: commercial / open source (GPLv2 with the (libstdc++)
              runtime exception), plus Freeware.
    \end{itemize}    
  }
\end{frame}

\begin{frame}
  \frametitle{Beneficios}
  \begin{itemize}
    \item \textbf{Abstracción de alto nivel}. \emph{TBB} se encarga de
          manejar los \emph{threads}, mientras que el usuario tan solo
          tiene que definir la \textbf{lógica de paralelismo}.
    \item \textbf{\emph{Threading} for performance}. \emph{TBB} está orientado
          a paralelizar tareas de computación intensiva.
    \item \textbf{Compatibilidad con otras librerías de paralelización}.
    \item \textbf{Data-parallel Programming}. Permite a diferentes
          \emph{threads} trabajar sobre partes del mismo conjunto de datos.
    \item \textbf{Generic Programming}. Los algoritmos son implementados
          sobre tipos genéricos, que el usuario especifica en función
          de su necesidad.
  \end{itemize}
\end{frame}


\begin{frame}
  \frametitle{Contenidos}
  Intel TBB proporciona:
  \begin{itemize}
  \item \textbf{Básicos:}\\\quad\texttt{parallel\_for}, \texttt{parallel\_reduce},
    \texttt{parallel\_scan}.
  \item \textbf{Avanzados:}\\\quad\texttt{parallel\_while}, \texttt{parallel\_do},
    \texttt{parallel\_sort}.
  \item \textbf{Contenedores:}\\\quad\texttt{concurrent\_queue},
    \texttt{concurrent\_vector}, \texttt{concurrent\_hash\_map}.
  \item \textbf{Exclusión mutua:}\\\quad\texttt{mutex},
    \texttt{spin\_mutex},\texttt{queuing\_mutex}.
  \item \textbf{Operaciones atómicas:}\\\quad
    \texttt{fetch\_and\_increment}, \texttt{fetch\_and\_decrement},
    \texttt{fecth\_and\_add}, \texttt{compare\_and\_swap},
    \texttt{compare\_and\_store}.
  \end{itemize}

\end{frame}



%===== EJEMPLO =====
\begin{frame}[fragile] % Frame ejemplo 1
  \frametitle{Ejemplo}
  \alert{Problema} de acceso a variable compartida:
  \begin{columns}
    \begin{column}{0.45\textwidth}
      \center Thread A:
      \begin{lstlisting}[language=C++, frame=single]
for( ; ; ){
    // do something ...
    // ...
    work_counter += n;
}
      \end{lstlisting}
    \end{column}
    \begin{column}{0.45\textwidth}
      \center Thread B:
      \begin{lstlisting}[language=C++, frame=single]
for( ; ; ){
    // do something ...
    // ...
    work_counter += n;
}
      \end{lstlisting}
    \end{column}
  \end{columns}
  \alert{Solución} usando pthreads:
  \begin{itemize}
  \item Proteger la variable (Mutex, cerrojos, semáforos, ...)
  \end{itemize}
\end{frame}

\begin{frame}[fragile] % Frame ejemplo 2
  \frametitle{Ejemplo}
  \alert{Solución} usando Intel TBB:

  \begin{lstlisting}[language=C++, frame=single]
tbb::atomic<int> work_counter;
// ...

for( ; ; ){
    // do something ...
    // ...
    work_counter += n;
}
  \end{lstlisting}
  Internamente está implementado:
  \begin{itemize}
  \item work\_counter.fetch\_and\_store
  \item work\_counter.fetch\_and\_add
  \item work\_counter.compare\_and\_swap
  \end{itemize}
\end{frame}

%======= BIBLIOGRAFÍA =======
\begin{frame}
  \frametitle{Bibliografía}
  
  \begin{enumerate}
  \item \textbf{Página oficial:}\\\url{https://www.threadingbuildingblocks.org/}
  \item \textbf{Intel TBB tutorial:}\\\url{http://www.inf.ed.ac.uk/teaching/courses/ppls/TBBtutorial.pdf}
  \item \textbf{Hpc wiki:}\\\url{https://wiki.hpcc.msu.edu/display/gpgpu/Intel+Threading+Building+Blocks+Talk+slides}
  \item \textbf{Wikipedia:}\\\url{https://en.wikipedia.org/wiki/Threading_Building_Blocks}
  \end{enumerate}
\end{frame}

\end{document}
